\documentclass[12pt]{article}
\usepackage[margin=0.915in]{geometry}

%%%
%%%             BEGIN Damian Preamble
%%%   Preface commands for pdfLaTeX output
%%%

%%
%% Calling relevant packages
%%

\usepackage[%
pdfstartview={FitV}, 
colorlinks=true,
menucolor=DarkGray,
linkcolor=MidnightBlue,
citecolor=MidnightBlue,
urlcolor=OrangeRed]{hyperref}
\usepackage{graphicx}
\DeclareGraphicsExtensions{.pdf}

%% From pandoc default: Begin
\usepackage{graphicx}
\makeatletter
\def\maxwidth{\ifdim\Gin@nat@width>\linewidth\linewidth\else\Gin@nat@width\fi}
\def\maxheight{\ifdim\Gin@nat@height>\textheight\textheight\else\Gin@nat@height\fi}
\makeatother
% Scale images if necessary, so that they will not overflow the page
% margins by default, and it is still possible to overwrite the defaults
% using explicit options in \includegraphics[width, height, ...]{}
\setkeys{Gin}{width=\maxwidth,height=\maxheight,keepaspectratio}
 %% From pandoc default: Begin

%\usepackage[hyper]{apacite}
\usepackage[svgnames]{xcolor}
\usepackage{rotating,bm,amsmath,amsfonts,amssymb,indentfirst,lscape,fancybox,fancyvrb,listings,pdfpages}
\usepackage[pagestyles]{titlesec}
%\usepackage{ucs}

%%
%% Format changes for chapter and section commands
%%

%\input{onesidedheader.tex}

%%
%% Chapter declarations
%%


%%
%% Section declarations
%%

%%
%% Subsection declarations
%%


%%
%% Subsubsection declarations
%%

%%
%% Listings setup information
%%

\lstset{language=R, frame=ltrb, framesep=5pt, xleftmargin=12pt, xrightmargin=5pt,
       numbers=none, breaklines=true, fancyvrb=true,
       breakatwhitespace=true, captionpos=b, abovecaptionskip=1.5ex,
       backgroundcolor=\color{Cornsilk},
       basicstyle=\small\color{DarkSlateGrey},
       keywordstyle=\ttfamily\color{DarkSlateGrey},
       identifierstyle=\ttfamily\color{DarkSlateGrey}\bfseries, 
       commentstyle=\color{Fuchsia},
       stringstyle=\ttfamily,
       showstringspaces=false}

%%
%% Header and Footer Specification 
%% TO ACTIVATE THE HEADER/FOOTER, ONE MUST PLACE \pagestyle{plain} 
%% followed by \pagestyle{damian} in the document
%%

\widenhead{0.14in}{0.14in}

\renewpagestyle{plain}{}

\newpagestyle{damian}[\sffamily]{
\headrule
\sethead[\sectiontitle][\chaptertitle][\thepage]
  {\sectiontitle}{\chaptertitle}{\thepage}

\footrule
\setfoot[][\raisebox{-.85ex}[0pt]{\NavigationBar}][]
{}{\raisebox{-.85ex}[0pt]{\NavigationBar}}{}
	\newcommand{\NavigationBar}{%
	  \Acrobatmenu{PrevPage}{Previous}\hspace{.5cm}
	  \Acrobatmenu{NextPage}{Next}\hspace{.5cm}
	  \Acrobatmenu{FirstPage}{First}\hspace{.5cm}
	  \Acrobatmenu{LastPage}{Last}\hspace{.5cm}
	  \Acrobatmenu{GoBack}{Back}\hspace{.5cm}
	  \Acrobatmenu{Quit}{Quit}%
}
}

%%
%% Custom specifications, commands, and colors
%%

\definecolor{DarkGray}{cmyk}{0,0,0,.624}
\newcommand{\R}{{\sffamily \textup{R}}}
\newcommand{\MlwiN}{{\sffamily \textup{MlwiN}}}
\newcommand{\Sweave}{{\sffamily \textup{Sweave}}}
\newcommand{\sssty}{\scriptscriptstyle}
\newcommand{\bigdot}{\ensuremath{\hspace{-.3ex}\bm{.}\hspace{-.05ex}}}
\renewcommand{\abstractname}{Abstract}
\renewcommand{\lstlistingname}{\R~Code Example}
\renewcommand{\arraystretch}{1.2}
\setlength{\fboxsep}{4mm}
\setlength{\fboxrule}{0.4pt}

%%
%%		END Damian Preamble
%%

\hypersetup{%
  pdftitle={ The Georgia Student Growth Model },
  pdfauthor={  Damian W. Betebenner  ,  Adam R. VanIwaarden  ,  \emph{National Center for the Improvement \ of Educational Assessment (NCIEA)}  },
  pdfcreator={ Damian W. Betebenner  , Adam R. VanIwaarden  , \emph{National Center for the Improvement \ of Educational Assessment (NCIEA)}  },
  pdfkeywords={},
  bookmarks=true} % pdfproducer={pdfLaTeX}

\usepackage{caption}
\usepackage{float}

\setcounter{secnumdepth}{3}

\usepackage{color}
\usepackage{fancyvrb}
\newcommand{\VerbBar}{|}
\newcommand{\VERB}{\Verb[commandchars=\\\{\}]}
\DefineVerbatimEnvironment{Highlighting}{Verbatim}{commandchars=\\\{\}}
% Add ',fontsize=\small' for more characters per line
\usepackage{framed}
\definecolor{shadecolor}{RGB}{248,248,248}
\newenvironment{Shaded}{\begin{snugshade}}{\end{snugshade}}
\newcommand{\KeywordTok}[1]{\textcolor[rgb]{0.13,0.29,0.53}{\textbf{{#1}}}}
\newcommand{\DataTypeTok}[1]{\textcolor[rgb]{0.13,0.29,0.53}{{#1}}}
\newcommand{\DecValTok}[1]{\textcolor[rgb]{0.00,0.00,0.81}{{#1}}}
\newcommand{\BaseNTok}[1]{\textcolor[rgb]{0.00,0.00,0.81}{{#1}}}
\newcommand{\FloatTok}[1]{\textcolor[rgb]{0.00,0.00,0.81}{{#1}}}
\newcommand{\CharTok}[1]{\textcolor[rgb]{0.31,0.60,0.02}{{#1}}}
\newcommand{\StringTok}[1]{\textcolor[rgb]{0.31,0.60,0.02}{{#1}}}
\newcommand{\CommentTok}[1]{\textcolor[rgb]{0.56,0.35,0.01}{\textit{{#1}}}}
\newcommand{\OtherTok}[1]{\textcolor[rgb]{0.56,0.35,0.01}{{#1}}}
\newcommand{\AlertTok}[1]{\textcolor[rgb]{0.94,0.16,0.16}{{#1}}}
\newcommand{\FunctionTok}[1]{\textcolor[rgb]{0.00,0.00,0.00}{{#1}}}
\newcommand{\RegionMarkerTok}[1]{{#1}}
\newcommand{\ErrorTok}[1]{\textbf{{#1}}}
\newcommand{\NormalTok}[1]{{#1}}



%\usepackage{pdfdraftcopy}
\newtheorem{proposition}{Proposition}
\newtheorem{theorem}{Theorem}
\newtheorem{definition}{Definition}
\newtheorem{corollary}{Corollary}
\DeclareMathOperator*{\argmin}{arg\,min}

% \usepackage{bbm}
\DeclareMathAlphabet{\mathbbm}{U}{bbm}{m}{n}
\SetMathAlphabet\mathbbm{bold}{U}{bbm}{bx}{n}
\DeclareMathAlphabet{\mathbbmss}{U}{bbmss}{m}{n}
\SetMathAlphabet\mathbbmss{bold}{U}{bbmss}{bx}{n}
\DeclareMathAlphabet{\mathbbmtt}{U}{bbmtt}{m}{n}

\begin{document}


\newcommand{\pl}[1]{\textsf{PL#1}}
\newcommand{\Cov}{\ensuremath{\mbox{\textsf{Cov}}}}
\newcommand{\Diag}{\ensuremath{\mbox{\textsf{Diag}}}}
\newcommand{\Bias}{\ensuremath{\mbox{\textsf{Bias}}}}
\newcommand{\Astar}[1]{\ensuremath{#1^{^*}}}
\thispagestyle{plain}
\pagestyle{damian}

\title{\textsf{\LARGE The Georgia Student Growth Model  \\\medskip A Technical Overview of the 2013-2014 Student Growth Percentile
Calculations }}
\author{  Damian W. Betebenner    \\   Adam R. VanIwaarden    \\   \emph{National Center for the Improvement \ of Educational Assessment (NCIEA)}   }

 \date{April 2015} 

\maketitle

\newpage

\begin{abstract}
DRAFT REPORT - DO NOT CITE! This report provides details about the
Georgia Student Growth Model methodology and presents a descriptive
analysis of the 2014 SGP calculation process and results.
\end{abstract}
\newpage

\section{Introduction}\label{introduction}

This report contains details on the 2013-2014 implementation of the
student growth percentiles (SGP) model for the state of Georgia. The
National Center for the Improvement of Educational Assessment (NCIEA)
contracted with the Georgia Department of Education (DOE) to implement
the SGP methodology using data derived from the Georgia Student
Assessment Program to create the
\href{http://www.gadoe.org/Curriculum-Instruction-and-Assessment/Assessment/Pages/Georgia-Student-Growth-Model.aspx}{Georgia
Student Growth Model}. The goal of the engagement with DOE is to create
a set of open source analytics techniques and conduct a set of initial
analyses that will eventually be conducted by DOE in following years.

The SGP methodology is an open source norm- and criterion-referenced
student growth analysis that produces student growth percentiles and
student growth projections/targets for each student with longitudinal
data in the state. The methodology is currently used for many purposes.
States and districts have used the results in various ways including
parent/student diagnostic reporting, institutional improvement, and
school and educator accountability. Specifics about the manner in which
growth is included in school and educator accountability can be found in
documents related to those accountability systems.

The report includes four sections covering Data, Analytics, SGP Results,
and Goodness of Fit:

\begin{itemize}
\itemsep1pt\parskip0pt\parsep0pt
\item
  \emph{Data} includes details on the decision rules used in the raw
  data preparation and student record validation.
\item
  \emph{Analytics} introduces some of the basic computational process
  and statistical methods implemented in 2013-2014.\footnote{More in
    depth treatment of the SGP Methodology can be found
    \href{https://github.com/CenterForAssessment/SGP_Resources/tree/master/articles}{here}.}
\item
  \emph{SGP Results} provides basic descriptive statistics from the 2014
  analyses.
\item
  \emph{Goodness of Fit} describes how well the statistical models used
  to produce SGPs fit Georgia students' data. This includes discussion
  of goodness of fit plots and the student- and school-level
  correlations between SGP and prior achievement.
\end{itemize}

Additionally, multiple appendices are included. Appendix A includes
Goodness of Fit plots for all content areas and grades. Appendix B
provides a technical description of the SIMEX correction for measurement
error with specific applications to Georgia.

\section{Data}\label{data}

Data for the Georgia Criterion-Referenced Competency Tests (CRCT) and
End of Course Tests (EOCT) used in the SGP analyses were supplied by the
Georgia DOE to the NCIEA for analysis in the fall of 2014. The current
longitudinal data set includes academic years 2006-2007 through
2013-2014. Subsequent years' analyses will augment this multi-year data
set allowing DOE to maintain a comprehensive longitudinal data set for
all students taking the CRCT and EOCT assessments and transitioning to
the
\href{http://www.gadoe.org/Curriculum-Instruction-and-Assessment/Assessment/Pages/Georgia-Milestones-Assessment-System.aspx}{Georgia
Milestones Assessment System.}

Student Growth Percentiles have been produced for students that have a
current score and at least one prior score in the same subject or a
related content area. SGPs were produced for grade-level (CRCT) Reading,
English Language Arts (ELA), Social Studies, Science and Mathematics.
For the 2014 academic year SGPs were produced for EOCT courses including
9\(^{th}\) Grade Literature, American Literature, U.S. History,
Economics, Biology, Physical Science, Mathematics II, Coordinate Algebra
and Analytic Geometry.

\subsection{Longitudinal Data}\label{longitudinal-data}

Growth analyses on assessment data require data that are linked to
individual students over time. Student growth percentile analyses
require, at a minimum two, and preferably three years of assessment data
for analysis of student progress. To this end it is necessary that a
unique student identifier be available so that student data records
across years can be merged with one another and subsequently examined.
Because some records in the assessment data set contain students with
more than one test score in a content area in a given year, a process to
create unique student records in each content area by year combination
was required in order to carry out subsequent growth analyses.
Furthermore, student records may be invalidated for other reasons. The
following business rules were used to either invalidate particular
student records or select the appropriate record for use in the
analyses.

\hyperdef{}{general-business-rules}{\subsubsection{General business
rules}\label{general-business-rules}}

\begin{enumerate}
\def\labelenumi{\arabic{enumi}.}
\itemsep1pt\parskip0pt\parsep0pt
\item
  Student records are invalidated if the student identifier is not
  exactly 10 digits long.
\item
  Student records with missing (``NA'') scores or scale scores outside
  of the possible range (usually 0) are invalidated.
\item
  Duplicate records that originated from the ``unmatched'' data files
  are invalidated.\footnote{Both ``Matched'' (where student to
    school/district links have been verified) and ``Unmatched'' data
    files are available. Unmatched data is added to the initial data
    file in hopes that it will possibly fill holes in students' score
    histories, but any current year SGPs produced will not be linked to
    any school, district or teacher.}
\item
  If a student has a main and a retest score\footnote{Retests were
    available in 2011-12 but retest data are no longer used in SGP
    calculations as of the 2012-13 school year.}, the record with the
  lower score is invalidated.
\item
  Student records with any administrative invalidation flag (for
  example, identifying test irregularities, students that did not
  attempt the test, or other issues) are invalidated.
\item
  School years for which CRCT tests were aligned to QCC curriculum are
  invalidated (SGPs were not produced in previous years and have not
  been used as prior scores subsequently).
\end{enumerate}

Beginning in 2014 the majority of the selection and invalidation of
student records was performed by GaDOE. Whereas in previous years NCIEA
staff conducted the data cleaning using \texttt{R} prior to running the
data analytics, GaDOE incorporated these business rules into the
\texttt{SQL} code used to pull data from GaDOE's data warehouse.

\subsubsection{CRCT specific business
rules}\label{crct-specific-business-rules}

\begin{enumerate}
\def\labelenumi{\arabic{enumi}.}
\itemsep1pt\parskip0pt\parsep0pt
\item
  If a student has multiple records (duplicate from the same subject,
  grade and administration period), their highest score was selected.
\item
  Student records from schools identified by the Governor's Office of
  Student Achievement (GOSA) as having artificially inflated student
  test scores from the 2008-2010 are invalidated.
\item
  If a student took more than one assessment in the same subject and
  school year but was identified as being in two different grades, the
  record with the highest grade level was selected.
\end{enumerate}

Table 1 shows the number of valid CRCT student records available for
analysis after applying the general and CRCT specific business
rules.\footnote{This number does not represent the number of SGPs
  produced, however, because students are required to have at least one
  prior score available as well.}

\begin{table}[H]
\caption*{\textbf{Table 1:} Number of Valid CRCT Student Records by Grade and Subject for 2014\label{table1}} 
\begin{center}
\begin{tabular}{lcllllll}
\hline\hline
\multicolumn{1}{c}{\bfseries }&\multicolumn{1}{c}{\bfseries }&\multicolumn{6}{c}{\bfseries Grades}\tabularnewline
\cline{1-8}
\multicolumn{1}{c}{Content Area}&\multicolumn{1}{c}{}&\multicolumn{1}{c}{3}&\multicolumn{1}{c}{4}&\multicolumn{1}{c}{5}&\multicolumn{1}{c}{6}&\multicolumn{1}{c}{7}&\multicolumn{1}{c}{8}\tabularnewline
\hline
Reading&&126,691&124,832&123,619&124,700&127,197&126,158\tabularnewline
ELA&&126,977&124,679&123,542&124,571&126,921&125,945\tabularnewline
Social Studies&&129,083&127,135&126,770&127,152&129,329&128,298\tabularnewline
Science&&129,604&127,612&127,219&127,540&129,772&128,864\tabularnewline
Mathematics&&127,463&124,665&122,976&124,027&126,346&125,192\tabularnewline
\hline
\end{tabular}\end{center}

\end{table}

\subsubsection{EOCT specific business
rules}\label{eoct-specific-business-rules}

\begin{enumerate}
\def\labelenumi{\arabic{enumi}.}
\itemsep1pt\parskip0pt\parsep0pt
\item
  If a student has multiple records (duplicate from the same subject and
  administration period), their highest score was selected.
\item
  For students who participate in two main administrations (i.e.,
  students who failed and retook an EOCT course), the first attempt is
  used as a prior to produce an SGP for the final attempt. In subsequent
  years, their final attempt may be used as a prior for other EOCT
  analyses.
\item
  Students who have tested out of EOCT courses are
  invalidated.\footnote{Beginning in the 2013-2014 school year, students
    had the opportunity to test out of an EOCT course by taking the test
    early and scoring at the Exceeds level. SGPs were not calculated for
    any test out attempt. Successful attempts, however, may be used as
    prior scores in subsequent years' analyses.}
\end{enumerate}

Table 2 shows the total number of valid EOCT student records available
for analysis after applying the general and EOCT specific business
rules.

\begin{table}[H]
\caption*{\textbf{Table 2:} Total Number of Valid EOCT Student Records by Subject for 2014\label{table2}} 
\begin{center}
\begin{tabular}{ll}
\hline\hline
\multicolumn{1}{c}{Content Area}&\multicolumn{1}{c}{Valid Records}\tabularnewline
\hline
Grade 9 Lit&134,863\tabularnewline
American Lit&114,511\tabularnewline
Economics&107,313\tabularnewline
US History&113,037\tabularnewline
Biology&131,978\tabularnewline
Physical Science&85,290\tabularnewline
Coordinate Algebra&143,211\tabularnewline
Analytic Geometry&120,224\tabularnewline
Mathematics II&10,432\tabularnewline
\hline
\end{tabular}\end{center}

\end{table}

\section{Analytics}\label{analytics}

This section provides basic details about the calculation of student
growth percentiles from Georgia state assessment data using the
\href{http://www.r-project.org/}{R Software Environment} in conjunction
with the \href{https://github.com/CenterForAssessment/SGP}{SGP Package}.
More in depth treatment of the data analysis process, including code
examples, has been provided to GaDOE staff in the 2014 Calculation
Manual.

Broadly, the SGP analysis of the Georgia longitudinal student assessment
data takes place in two steps:

\begin{enumerate}
\def\labelenumi{\arabic{enumi}.}
\itemsep1pt\parskip0pt\parsep0pt
\item
  Data Preparation
\item
  Data Analysis
\end{enumerate}

Those familiar with data analysis know that the bulk of the effort in
the above two step process lies with Step 1: Data Preparation. Following
thorough data cleaning and preparation, data analysis using the SGP
Package takes clean data and makes it as easy as possible to calculate,
summarize, output and visualize the results from SGP analyses.

\subsection{Data Preparation}\label{data-preparation}

The data preparation step involves taking data provided by the GaDOE and
producing a \texttt{.Rdata} file that will subsequently be analyzed in
Step 2. This process is carried out annually as new data becomes
available from the state assessment program. The data supplied by the
GaDOE Information Technology department and subsequently cleaned and
processed using a two step process:

\textbf{Step 1a.} Initial data extraction and cleaning

In this first step a clean data set is extracted from the Georgia
student data warehouse using an internal \texttt{SQL} connection and
command script. Through this process, student records are selected and
invalidated based upon the \hyperref[general-business-rules]{business
rules discussed above}. The end result is an exported, pipe-delimited
file where a each valid student record is unique by content area, school
year, student identifier (GTID), and test administration period.

\textbf{Step 1b.} Final data cleaning and preparation in \texttt{R}

In this step the clean data from step 1a is read into
\href{http://www.r-project.org/}{\texttt{R}} and modified slightly with
regard to variable types/classes. The result is an \texttt{.Rdata} file
that is suitable for analysis with the
\href{https://github.com/CenterForAssessment/SGP}{SGP Package}. In
previous years' analyses the bulk of the data cleaning and
implementation of business rule validation from step 1a was also
performed in \texttt{R} by NCIEA staff. However, the division of the
cleaning and validation tasks used in 2014 is an important step in
moving towards GaDOE self-sufficiency in calculating growth percentiles
in subsequent years.

The cleaned and formatted data is finally combined with the existing
data housed in the Georgia
\href{http://cran.r-project.org/web/packages/SGP/SGP.pdf}{SGP class
object}, which was created through the 2013 analyses. With an
appropriate longitudinal data prepared, we move to the calculation and
estimation of student-level SGPs.

\subsection{2014 Data Analysis}\label{data-analysis}

The objective of the student growth percentile (SGP) analysis is to
describe how (a)typical a student's growth is by examining his/her
current achievement relative to students with a similar achievement
history; i.e his/her \emph{academic peers} (see
\href{https://docs.google.com/viewer?url=http\%3A\%2F\%2Fwww.gadoe.org\%2FCurriculum-Instruction-and-Assessment\%2FAssessment\%2FDocuments\%2FSGP\%2520FAQ\%2520072214.pdf}{Section
2 of the GSGM FAQ}). The estimation of this norm-referenced growth
quantity is conducted using quantile regression (Koenker, 2005) to model
curvilinear functional relationships between student's prior and current
scores. One hundred such regression calculations are run for each
separate analysis (defined as a unique year by content area by grade by
prior order combination). The end product of these 100 separate
regression models is a single coefficient matrix, which serves as a
look-up table to relate prior student achievement to current achievement
for each percentile. This process ultimately leads to the calculation of
tens of thousands of calculations per year (and many more when SIMEX
measurement error corrections are performed) during each of Georgia's
annual batch of analyses. For a more in-depth discussion of the
calculation and estimation of SGPs, see Betebenner ((2009)), and see
Shang, VanIwaarden and Betebenner ((Shang, VanIwaarden, \& Betebenner,
2015)) and Appendix B of this report for further information on the
SIMEX measurement error correction methodology.

The 2014 Georgia SGP analyses follow a work flow established in previous
years that includes the following 6 steps:

\begin{enumerate}
\def\labelenumi{\arabic{enumi}.}
\itemsep1pt\parskip0pt\parsep0pt
\item
  Update the Georgia assessment meta-data required for SGP calculations
  using the \texttt{SGP} Package.
\item
  Create annual SGP configurations for EOCT analyses.
\item
  Create any baseline matrices and SIMEX baseline coefficient matrices
  needed for new content areas sequences.
\item
  Conduct CRCT SGP Analyses.
\item
  Conduct EOCT SGP Analyses.
\item
  Combine results into the master longitudinal data set, summarize
  results and output data.
\end{enumerate}

\subsubsection{Update Georgia meta-data in \texttt{SGP}
Package}\label{update-georgia-meta-data-in-sgp-package}

The use of higher-level functions included in the SGP package (e.g.
\texttt{analyzeSGP}) requires the availability of state specific
assessment information. This meta-data is compiled in a \texttt{R}
object named \texttt{SGPstateData} that is housed in the package. The
required updates for the 2014 analyses included a) the additions of
knots and boundaries, proficiency level cutscores, and conditional
standard errors of measurement (CSEMs) for Analytic Geometry, b) adding
a new variable to the Variable Name Lookup table, and c) updating the
norm group preferences object.

\textbf{Calculation and addition of knots and boundaries}

Calculation of SGPs includes the use of cubic B-spline basis functions
to more adequately model the heteroscedasticity and non-linearity found
in assessment data.\footnote{It should be noted that the independent
  estimation of the regression functions can potentially result in the
  crossing of the quantile functions. This occurs near the extremes of
  the distributions and is potentially more likely to occur given the
  use of non-linear functions. A potential result of allowing the
  quantile functions to cross would be \emph{lower} percentile
  estimations of growth for \emph{higher} observed scale scores at the
  extremes (give all else equal in prior scores) and vice versa. In
  order to deal with these contradictory estimates, quantile regression
  results are isotonized to prevent quantile crossing following the
  methods derived by Chernozhukov, Fernandez-Val and Glichon (2010).}
These functions require the selection of boundary and interior knots.
Boundary knots are end-points outside of the scale score distribution
that anchor the B-spline basis. These are generally selected by
extending the entire range of scale scores by 10\%. That is, they are
defined as lying 10\% below the lowest obtainable (or observed) scale
score (LOSS) and 10\% above the highest obtainable scale score (HOSS).
The interior knots are the \emph{internal} breakpoints that define the
spline. The default choice in the \texttt{SGP} package (Betebenner,
{VanIwaarden}, Domingue, \& Shang, 2015) is to select the 20\(^{th}\),
40\(^{th}\), 60\(^{th}\) and 80\(^{th}\) quantiles of the observed scale
score distribution.

In general the knots and boundaries are computed from a distribution
comprised of several years of test data (i.e.~multiple cohorts) so that
any irregularities in a single year are smoothed out. This is important
because subsequent annual analyses use these same knots and boundaries
as well. All defaults were used to compile the knots and boundaries for
Georgia from the CRCT and EOCT tests in previous years, and were also
used in 2014 to compute the Analytic Geometry knots and boundaries. New
knots and boundaries will be required for Georgia Milestones assessments
beginning in 2016 at which point they will be used as the dependent
variables in the quantile regressions.

\textbf{Proficiency level cutscores}

Cutscores, which are set externally by the GaDOE through
standard-setting processes, are mainly required for student growth
projections, which were not computed as part of the 2014
analyses.\footnote{Projections were not included in the 2014 analyses
  due to the switch to Georgia Milestones Assessments in 2015. Student
  Growth Projections assume consistency in assessment programs. It would
  therefore would be nonsensical to project progress toward tests that
  will never be taken.} However, they may likely be used in future years
and so were added at this point.

\textbf{Conditional standard errors of measurement (CSEMs)}

The calculation of SIMEX adjusted student growth percentiles requires
the availability of each assessments' standard errors of measurement.
The CSEM data for all other content areas had been compiled and added in
previous years, but were required for Analytic Geometry in this initial
year of testing. Raw CSEM data was provided by GaDOE in spreadsheets and
the NCIEA compiled this into an appropriate \texttt{R} data object for
inclusion in the \texttt{SGPstateData} file.

\textbf{The ``Variable Name Lookup'' table}

The \texttt{SGPstateData} file also includes a lookup table that allows
various functions in the SGP package to translate between the naming
conventions used within them and the variable names that the GaDOE uses.
The 2014 data included a gifted and talented student identifier, which
was used in result summarizations.

\textbf{Norm group preferences}

The process through which EOCT analyses are run can produce multiple
SGPs for some students. In order to identify which quantity will be used
as the students' ``official'' SGP and subsequently merged into the
master longitudinal data set, a system of norm group preferencing is
established and is encoded into a lookup table and included in the
\texttt{SGPstateData}. In general, the preference is given to:

\begin{itemize}
\itemsep1pt\parskip0pt\parsep0pt
\item
  Progressions with the greatest number of prior scale scores.
\item
  Progressions in which a student has repeated a course.
\item
  Progressions that do not include a skipped year (i.e.~a gap in the
  scale score history).
\item
  Progressions for block-schedule course taking patterns.
\end{itemize}

The next section describes the process by which the individual course
progression analyses are established and how the preferencing system is
included within their configuration code.

\subsubsection{Create annual SGP
configurations.}\label{create-annual-sgp-configurations.}

Unlike CRCT analyses, EOCT analyses are specialized enough so that it is
necessary to specify the analyses to be performed via explicit
configuration code. For several years, configurations have been employed
to conduct EOCT SGP analyses for Georgia. The configurations are broken
up into four separate R scripts based on content domain (literature,
mathematics, science and social studies). Each configuration specifies a
set of parameters that defines the norm group of students to be
examined. Every potential norm group is defined by, at a minimum, the
progressions of content area, academic year and grade-level. Each
configuration used for the Georgia EOCT analyses contain these elements:

\begin{itemize}
\itemsep1pt\parskip0pt\parsep0pt
\item
  \textbf{\texttt{sgp.content.areas}:} A progression of values that
  specifies the content areas to be looked at and the order in which the
  courses were taken.
\item
  \textbf{\texttt{sgp.panel.years}:} The progression of the years
  associated with the content area progression
  (\texttt{sgp.content.areas}) provided in the configuration,
  potentially allowing for skipped years, block schedules, etc.
\item
  \textbf{\texttt{sgp.grade.sequences}:} The grade progression
  associated with the content area and year progressions provided in the
  configuration. \emph{`EOCT'} stands for `End Of Course Test'. The use
  of the generic \emph{`EOCT'} allows for secondary students to be
  compared based on the pattern of course taking rather than being
  dependent upon grade-level/class-designation.
\item
  \textbf{\texttt{sgp.panel.years.within}:} A vector of same length as
  the year progression (sgp.panel.years) indicating what observation is
  to be used for the individual student (when multiple observations
  exist within a single year). Typically the ``Last'' observation is
  used as the prior score (covariate) and the ``First'' observation is
  used as the current year score (outcome).
\item
  \textbf{\texttt{sgp.exact.grade.progression}:} A Boolean argument (set
  to TRUE) indicating whether to run the EXACT configuration as written
  (rather than taking progressively restricted nested subsets of the
  configuration if FALSE).
\item
  \textbf{\texttt{sgp.calculate.simex}:} A Boolean argument indicating
  whether cohort-referenced SIMEX adjustment analyses should be run as
  part of the analysis for this configuration. Excluding the argument
  (or explicitly setting to NULL) has the same effect as setting it to
  FALSE.
\item
  \textbf{\texttt{sgp.calculate.simex.baseline}:} A Boolean argument
  indicating whether baseline-referenced SIMEX adjustment analyses
  should be run as part of the analysis for this configuration.
  Excluding the argument (or explicitly setting to NULL) has the same
  effect as setting it to FALSE.
\item
  \textbf{\texttt{sgp.norm.group.preference}:} Because a student can
  potentially be included in more than one analysis/configuration, this
  argument provides a ranking specifying which SGP is preferable for
  being matched with the student in the
  \href{https://github.com/CenterForAssessment/SGP/blob/master/R/combineSGP.R}{\texttt{combineSGP}}
  step. \emph{Lower numbers correspond with higher preference.}
\end{itemize}

As an example, here is one of the Biology configurations used to defined
the 2014 SGP analyses:

\begin{Shaded}
\begin{Highlighting}[]
\NormalTok{...}

  \NormalTok{BIOLOGY}\FloatTok{.2014} \NormalTok{=}\StringTok{ }\KeywordTok{list}\NormalTok{(}
        \DataTypeTok{sgp.content.areas=}\KeywordTok{c}\NormalTok{(}\StringTok{'SCIENCE'}\NormalTok{, }\StringTok{'PHYSICAL_SCIENCE'}\NormalTok{, }\StringTok{'BIOLOGY'}\NormalTok{),}
        \DataTypeTok{sgp.panel.years=}\KeywordTok{c}\NormalTok{(}\StringTok{'2012'}\NormalTok{, }\StringTok{'2013'}\NormalTok{, }\StringTok{'2014'}\NormalTok{),}
        \DataTypeTok{sgp.grade.sequences=}\KeywordTok{list}\NormalTok{(}\KeywordTok{c}\NormalTok{(}\StringTok{'8'}\NormalTok{, }\StringTok{'EOCT'}\NormalTok{, }\StringTok{'EOCT'}\NormalTok{)),}
        \DataTypeTok{sgp.panel.years.within=}\KeywordTok{c}\NormalTok{(}\StringTok{'LAST_OBSERVATION'}\NormalTok{, }\StringTok{'LAST_OBSERVATION'}\NormalTok{, }\StringTok{'FIRST_OBSERVATION'}\NormalTok{),}
        \DataTypeTok{sgp.exact.grade.progression=}\OtherTok{TRUE}\NormalTok{,}
        \DataTypeTok{sgp.calculate.simex=}\OtherTok{TRUE}\NormalTok{,}
        \DataTypeTok{sgp.calculate.simex.baseline=}\OtherTok{TRUE}\NormalTok{,}
        \DataTypeTok{sgp.norm.group.preference=}\DecValTok{4}\NormalTok{), }
\NormalTok{...}
\end{Highlighting}
\end{Shaded}

Configurations are \texttt{R} scripts that are sourced as part of the
larger SGP analysis to be discussed later and to construct the norm
group preference object discussed previously.

\subsubsection{Create baseline and SIMEX baseline coefficient
matrices.}\label{create-baseline-and-simex-baseline-coefficient-matrices.}

The 2014 CRCT \& EOCT analyses Georgia employed baseline-referenced and
SIMEX adjusted baseline-referenced SGPs. For most grade and content area
analyses the coefficient matrices required to produce these results were
produced prior to 2014. However, baseline matrices for several EOCT
course progressions were found to have adequate data for baseline matrix
construction.

Baseline coefficient matrices differ form traditional cohort-referenced
matrices in that they are produced using a ``super-cohort'' of students
rather than a single year cohort. For example, the baseline super-cohort
for the Biology course sequence shown above consisted of three combined
cohorts using GPS Science and EOCT Physical Science and Biology test
scores from 2008 to 2012. Every student who took Biology in the 2010,
2011 and 2012 academic years with consecutive Physical Science and
8\(^{th}\) grade Science prior scale scores would have been included in
the calculations of the Growth Percentile regression equations. Because
most baseline coefficient matrices were produced using previous years'
data, the 2014 cohort data does not contribute to the calculation of the
models used to produce their SGPs. Additional general information about
baseline analyses can be found in the discussion of
\hyperref[baseline-referenced-median-sgps]{Median SGP results below}.

The baseline matrices produced in 2014 were created using the same years
of data as those previously produced in those content areas. American
Literature and 9\(^{th}\) Grade Literature matrices (adding in skipped
year patterns not previously modeled) were produced using data from 2007
to 2012 and US History matrices were produced using data from 2008 to
2013. These matrices were then added to the existing collection and
their course progressions added in the 2014 configuration scripts.

The baseline-referenced, SIMEX-corrected versions of these matrices were
then produced using the uncorrected versions as the ``naive'' model.
Details of the SIMEX procedure can be found in Appendix B of this
report. In all 200 - 400 matrices were produced for each course
progression depending on the number of priors used.

\subsubsection{Conduct CRCT SGP
Analyses}\label{conduct-crct-sgp-analyses}

The CRCT grade level ELA, Reading, Math, Science and Social Studies
analyses were conducted separately from the EOCT course analyses. Mainly
this is done because the CRCT analyses do not require custom
configuration scripts like EOCT does. Given the simplistic use of only
the same test as prior test scores in these analyses, the SGP package's
functions can identify the analyses to be conducted internally. This
also allows the analyses to be completed and checked, and interim
results saved which is sometimes necessary due to the computationally
intensive nature of the analysis, particularly with a state the size of
Georgia.

The \texttt{analyzeSGP} function was used to calculate all SGPs. Initial
2014 analyses included the baseline and SIMEX corrected baseline SGP
production (the ``official'' SGPs), as well as annual cohort-referenced
SGPs.\footnote{Cohort referenced SGPs are not used in a official
  capacity, but are calculated as an additional resource for GaDOE. Note
  that CRCT cohort-referenced SGPs have used all available prior
  information, whereas Baseline referenced SGPs use only 2 prior years
  of data.} SIMEX corrected cohort-referenced SGPs were produced
subsequent to requests made by the Georgia Educator Effectiveness
Technical Advisory Committee.

\subsubsection{Conduct EOCT SGP
Analyses}\label{conduct-eoct-sgp-analyses}

The EOCT content area progression specific analyses were then calculated
using the configuration scripts described above. The baseline-referenced
content areas were analyzed separately from the cohort-referenced
subjects (mathematics domain). This was done due to the computationally
intensive nature of the calculations rather than necessity, as all
relevant argument controls are specified in the configuration scripts.
As with the CRCT analyses, cohort-referenced SIMEX corrected versions of
all SGPs were eventually produced.\footnote{Unlike the CRCT, EOCT
  cohort-referenced SGPs use only 2 prior years of data as specified in
  the configuration scripts.}

\subsubsection{Merge 2014 results into the longitudinal data, summarize
and
output.}\label{merge-2014-results-into-the-longitudinal-data-summarize-and-output.}

Once all analyses are completed the results are merged into the master
longitudinal data set in the \texttt{@Data} slot of the SGP class object
using the \texttt{combineSGP} function. The data is then summarized
using the \texttt{summarizeSGP} function, which produced many tables of
descriptive statistics that are disaggregated at the state, district,
school and instructor levels. These basic summary tables are also
further disaggregated by the demographic groups available in the data
set and listed in
\href{https://github.com/CenterForAssessment/SGPstateData/blob/master/Variable_Name_Lookup/GA_Variable_Name_Lookup.csv}{Georgia's
Variable Name Lookup table}. Finally, a pipe delimited version of the
complete long data is output. Additionally for Georgia, the NCIEA
produces a formatted version of the 2014 results, which contains fields
needed for rendering data in the state's
\href{http://gastudentgrowth.gadoe.org/}{visualization tool} such as
students' entire prior score and course progression history.

\section{SGP Results}\label{sgp-results}

In the following sections basic descriptive statistics from the 2014
analyses are provided, including the state-level mean and median growth
percentiles. Currently Georgia uses either cohort or baseline-referenced
SGPs as the official student-level growth metric, depending upon the
content area. The goal is to eventually use only baseline-referenced
SGPs, however this requires an assessment program to have been
established for several years. Given that EOCT math courses have
undergone changes in recent years, only cohort-referenced SGPs are
currently available for them. Similarly, the upcoming transition from
the CRCT and EOCT to
\href{http://www.gadoe.org/Curriculum-Instruction-and-Assessment/Assessment/Pages/Georgia-Milestones-Assessment-System.aspx}{Georgia
Milestones} assessments will require the exclusive use of
cohort-referenced analyses for several years until baseline referencing
is again feasible.

Georgia DOE has also recently decided to implement the SIMEX measurement
error correction method in the calculation of student-level SGPs. This
correction has been applied to both cohort and baseline-referenced SGPs
and the descriptive statistics from these analyses are also presented
here. The interested reader can find more in depth discussions of the
cohort and baseline-referenced SGP methodology in the
\href{https://github.com/CenterForAssessment/SGP_Resources/tree/master/articles}{available
literature}, and information about the SIMEX measurement error
correction methodology is available in Appendix B of this report as well
as in recent academic articles (Shang et al., 2015).

\subsection{Cohort Referenced Median
SGPs}\label{cohort-referenced-median-sgps}

Growth percentiles, being quantities associated with each individual
student, can be easily summarized across numerous grouping indicators to
provide summary results regarding growth. The median and mean of a
collection of growth percentiles are used as measures of central
tendency that summarize the distribution as a single number. With
perfect data fit, we expect the state median of all student growth
percentiles in any grade to be 50 because the data are norm-referenced
across all students in the state. Median (and mean) growth percentiles
well below 50 represent growth less than the state ``average'' and
median growth percentiles well above 50 represent growth in excess of
the state ``average''.

To demonstrate the norm-referenced nature of the growth percentiles
viewed at the state level, Table 3 presents the cohort-referenced growth
percentile medians and means for the EOCT mathematics based content
areas.

\begin{table}[H]
\caption*{\textbf{Table 3:} EOCT Median and Mean \emph{Cohort} Referenced Student Growth Percentile by Content Area for 2014\label{table3}} 
\begin{center}
\begin{tabular}{lll}
\hline\hline
\multicolumn{1}{c}{Content Area}&\multicolumn{1}{c}{Median SGP}&\multicolumn{1}{c}{Mean SGP}\tabularnewline
\hline
Coordinate Algebra&50&49.7\tabularnewline
Analytic Geometry&49&49.4\tabularnewline
Mathematics II&50&49.8\tabularnewline
\hline
\end{tabular}\end{center}

\end{table}

Based upon perfect model fit to the data, the median of all state growth
percentiles in each grade by year by subject combination should be 50.
That is, in the conditional distributions, 50 percent of growth
percentiles should be less than 50 and 50 percent should be greater than
50. Deviations from 50 indicate imperfect model fit to the data.
Imperfect model fit can occur for a number of reasons, some due to
issues with the data (e.g., floor and ceiling effects leading to a
``bunching'' up of the data) as well as issues due to the way that the
SGP function fits the data. The results in Table 3 are close to perfect,
with almost all values equal to 50.

The results are coarse in that they are aggregated across tens of
thousands of students. More refined fit analyses are presented in the
\hyperref[goodness-of-fit]{Goodness-of-Fit section that follows}.
Depending upon feedback from Georgia DOE, it may be desirable to tweak
with some operational parameters and attempt to improve fit even
further. The impact upon the operational results based on better fit is
expected to be extremely minor.

It is important to note how, at the entire state level, the
\emph{norm-referenced} growth information returns little information on
annual trends due to its norm-reference nature. What the results
indicate is that a typical (or average) student in the state
demonstrates 50\(^{th}\) percentile growth. That is, ``typical
students'' demonstrate ``typical growth''. One benefit of the
norm-referenced results follows when subgroups are examined (e.g.,
schools, district, demographic groups, etc.) Examining subgroups in
terms of the median of their student growth percentiles, it is then
possible to investigate why some subgroups display lower/higher student
growth than others. Moreover, because the subgroup summary statistic
(i.e., the median) is composed of many individual student growth
percentiles, one can break out the result and further examine the
distribution of individual results.

\hyperdef{}{baseline-referenced-median-sgps}{\subsection{Baseline
Referenced Median SGPs}\label{baseline-referenced-median-sgps}}

Baseline SGPs allow one to look at norm-referenced growth through
another lens. Rather than considering a single year's cohort, baseline
SGPs are referenced against a ``super-cohort'' consisting of several
years of students linked by common course/grade progressions. This
super-cohort serves as a baseline against which future annual cohorts
can be compared, which allows for the examination of whether or not the
system as a whole might be improving (or declining) in terms of growth
over time relative to the baseline group. If the system is improving
over time, we would see those higher levels of effectiveness in the form
of median SGPs that are greater than 50 (what \emph{was} a higher rate
of growth in the past is now typical growth). An assumption required for
the use of baseline-referenced SGPs is that the scale scores are well
anchored. If this assumption is violated, then any deviation from
``typical'' growth may be purely an artifact of the test scaling
procedure. States transitioning from old test to new tests, for example,
must reset their baseline growth because a new scale is in effect. Table
4 provides the median and mean baseline-referenced results from
Georgia's CRCT assessments and Table 5 provides the EOCT results.

\begin{table}[H]
\caption*{\textbf{Table 4:} CRCT Median (Mean) \emph{Baseline} Student Growth Percentile by Grade and Content Area for 2014\label{}} 
\begin{center}
\begin{tabular}{lclllll}
\hline\hline
\multicolumn{1}{c}{\bfseries }&\multicolumn{1}{c}{\bfseries }&\multicolumn{5}{c}{\bfseries Grades}\tabularnewline
\cline{1-7}
\multicolumn{1}{c}{Content Area}&\multicolumn{1}{c}{}&\multicolumn{1}{c}{4}&\multicolumn{1}{c}{5}&\multicolumn{1}{c}{6}&\multicolumn{1}{c}{7}&\multicolumn{1}{c}{8}\tabularnewline
\hline
Reading&&59 (56.6)&48 (48.7)&61 (58)&58 (55.7)&60 (55.9)\tabularnewline
ELA&&48 (48.3)&46 (47.1)&43 (45.2)&57 (54.8)&45 (46.1)\tabularnewline
Social Studies&&50 (49.5)&48 (48.3)&52 (51.5)&48 (48.5)&47 (47.8)\tabularnewline
Science&&55 (53.2)&50 (49.7)&46 (47.5)&54 (52.8)&54 (52.4)\tabularnewline
Mathematics&&49 (49.6)&44 (45.1)&56 (54.1)&37 (40.6)&52 (50.4)\tabularnewline
\hline
\end{tabular}\end{center}

\end{table}

\begin{table}[H]
\caption*{\textbf{Table 5:} EOCT Median and Mean \emph{Baseline} Referenced Student Growth Percentile by Content Area for 2014\label{}} 
\begin{center}
\begin{tabular}{lll}
\hline\hline
\multicolumn{1}{c}{Content Area}&\multicolumn{1}{c}{Median SGP}&\multicolumn{1}{c}{Mean SGP}\tabularnewline
\hline
Grade 9 Lit&51&50.7\tabularnewline
American Lit&48&48.7\tabularnewline
Economics&52&51.0\tabularnewline
US History&45&46.5\tabularnewline
Biology&43&45.0\tabularnewline
Physical Science&48&48.3\tabularnewline
\hline
\end{tabular}\end{center}

\end{table}

\subsection{SIMEX Adjusted Student Growth
Percentiles}\label{simex-adjusted-student-growth-percentiles}

The use of error-prone standardized test scores in statistical models
can lead to numerous problems. Understanding the effects of measurement
error and correcting it is particularly difficult in the SGP model given
that it is based on non-parametric quantile regression. Preliminary
investigations suggest that the use of error-prone measures may lead SGP
estimates to be inflated for students with high prior achievement and
underestimated for students with lower prior achievement. This bias at
the individual student-level can translate to bias in aggregate measures
(such as median and mean SGPs) when student sorting based on prior
achievement exists at the level of aggregation (e.g. classrooms or
schools). As a result, growth and prior achievement are (positively)
correlated, giving schools and teachers with higher achieving students
an undue advantage and disadvantaging those with a preponderance of low
achieving students.

Simulation-extrapolation (SIMEX) methods of correcting for measurement
error induced bias in the SGP estimation has been proposed and recently
tested in Georgia. Descriptive statistics from these analyses are
provided here and in the \hyperref[goodness-of-fit]{Goodness-of-Fit}
section. Additional technical information about the SIMEX procedure in
general and its use in the calculation of cohort and baseline-referenced
SGPs is provided in Appendix B of this report.

\begin{table}[H]
\caption*{\textbf{Table 6:} SIMEX Corrected EOCT Median and Mean \emph{Cohort} Referenced Student Growth Percentile by Content Area for 2014\label{table6}} 
\begin{center}
\begin{tabular}{lll}
\hline\hline
\multicolumn{1}{c}{Content Area}&\multicolumn{1}{c}{Median SGP}&\multicolumn{1}{c}{Mean SGP}\tabularnewline
\hline
Coordinate Algebra&50&50.1\tabularnewline
Analytic Geometry&50&49.7\tabularnewline
Mathematics II&50&50.1\tabularnewline
\hline
\end{tabular}\end{center}

\end{table}

\begin{table}[H]
\caption*{\textbf{Table 7:} SIMEX Corrected CRCT Median (Mean) \emph{Baseline} Student Growth Percentile by Grade and Content Area for 2014\label{table7}} 
\begin{center}
\begin{tabular}{lclllll}
\hline\hline
\multicolumn{1}{c}{\bfseries }&\multicolumn{1}{c}{\bfseries }&\multicolumn{5}{c}{\bfseries Grades}\tabularnewline
\cline{1-7}
\multicolumn{1}{c}{Content Area}&\multicolumn{1}{c}{}&\multicolumn{1}{c}{4}&\multicolumn{1}{c}{5}&\multicolumn{1}{c}{6}&\multicolumn{1}{c}{7}&\multicolumn{1}{c}{8}\tabularnewline
\hline
Reading&&57 (55.9)&47 (47.9)&60 (57.5)&58 (54.7)&60 (55.2)\tabularnewline
ELA&&46 (48)&46 (46.6)&42 (44.6)&57 (54.8)&44 (45.7)\tabularnewline
Social Studies&&48 (49.2)&47 (48.2)&52 (51.7)&48 (48.2)&47 (47.7)\tabularnewline
Science&&54 (53)&50 (49.4)&46 (47.3)&54 (52.8)&53 (52.1)\tabularnewline
Mathematics&&50 (49.3)&43 (44.7)&56 (54.1)&36 (39.9)&51 (50.3)\tabularnewline
\hline
\end{tabular}\end{center}

\end{table}

\begin{table}[H]
\caption*{\textbf{Table 8:} SIMEX Corrected EOCT Median and Mean \emph{Baseline} Referenced Student Growth Percentile by Content Area for 2014\label{table8}} 
\begin{center}
\begin{tabular}{lll}
\hline\hline
\multicolumn{1}{c}{Content Area}&\multicolumn{1}{c}{Median SGP}&\multicolumn{1}{c}{Mean SGP}\tabularnewline
\hline
Grade 9 Lit&50&50.2\tabularnewline
American Lit&47&47.9\tabularnewline
Economics&50&50.5\tabularnewline
US History&43&45.9\tabularnewline
Biology&42&44.5\tabularnewline
Physical Science&47&47.9\tabularnewline
\hline
\end{tabular}\end{center}

\end{table}

A comparison of the unadjusted (Tables 3 to 5) and SIMEX corrected
(Tables 6 through 8) shows very little difference in the medians and
means. This is not surprising as the majority of the growth percentiles
for students in the middle of the prior score distributions change very
little, and the larger changes that occur for students in the extremes
of the prior score distributions tend to even out.

\hyperdef{}{goodness-of-fit}{\section{Goodness of
Fit}\label{goodness-of-fit}}

Examination of goodness-of-fit was conducted by comparing the estimated
conditional density with the theoretical uniform density of the SGPs.
Despite the use of B-splines to accommodate heteroscedasticity and
skewness of the conditional density, assumptions are made concerning the
number and position of spline knots that impact the percentile curves
that are fit. With an infinite population of test takers, at each prior
scaled score, with perfect model fit, the expectation is to have 10
percent of the estimated growth percentiles between 1 and 9, 10 and 19,
20 and 29, \ldots{}, and 90 and 99. Deviations from 10 percent would be
indicative of lack of model fit.

\subsection{Model Fit Plots}\label{model-fit-plots}

Using all available CRCT and EOCT scores as the variables, estimation of
student growth percentiles was conducted for each possible student
(those with a current score and at least one prior score). Percentages
of student growth percentiles between the 10\(^{th}\), 20\(^{th}\),
30\(^{th}\), 40\(^{th}\), 50\(^{th}\), 60\(^{th}\), 70\(^{th}\),
80\(^{th}\), and 90\(^{th}\) percentiles were calculated based upon the
decile of the prior year's scaled score (the total students in each for
the analyses varies depending on grade and subject). Results for the
B-spline parameterizations for all CRCT and EOCT analyses are given in
Appendix A.

\subsubsection{Cohort referenced model
fit}\label{cohort-referenced-model-fit}

The results in all subjects are excellent with a few exceptions.
Deviations from perfect fit are indicated by red and blue shading. The
further \emph{above} 10 the darker the red, and the further \emph{below}
10 the darker the blue. In instances where large deviations from 10
occur, the likely cause is that there is a mass point associated with
certain scale scores that makes it impossible to ``split'' the score at
a dividing point forcing a majority of the scores into an adjacent cell.
This is the case with all large deviations observed in the Georgia data.

Figure 1 shows the results for 8\(^{th}\) grade ELA. Although this CRCT
subject is officially a baseline-referenced subject, we provide the plot
here as an exemplar of cohort-referenced model fit with which to compare
with baseline and SIMEX model fit using the same student data in
subsequent sections.

\begin{figure}[htbp]
\centering
\includegraphics{img/Examples/Goodness_of_Fit/ELA.2014/gofSGP_Grade_8.png}
\caption{Goodness of Fit Plot for 2014 \textbf{\emph{Cohort}} Referenced
8\(^{th}\) Grade ELA.}
\end{figure}

\subsubsection{Baseline referenced model
fit}\label{baseline-referenced-model-fit}

Assumptions regarding model fit are much different for
baseline-referenced models. Perfect model fit is not expected, or even
desired, in these analyses because the data from the 2014 cohort was not
used to produce the baseline model parameter estimates. As discussed
above, these models may allow us to see whether or not the system as a
whole might be improving (or declining) over time. These changes would
be indicated by model \emph{misfit} in the goodness of fit plots. A
greater number of students with high growth (red cells on the right hand
side of the first plots) would suggest system improvement. Ideally we
would see this for students in the entire range prior achievement.

Figure 2 shows how the baseline model fits the 2014 8\(^{th}\) grade ELA
student data. Here we see that the model suggests student performance
has declined relative to the established baseline. This appears to be
generally true for students from all prior achievement levels with the
exception of students in the lowest decile of prior scores.

\begin{figure}[htbp]
\centering
\includegraphics{img/Examples/Goodness_of_Fit/ELA.2014.BASELINE/gofSGP_Grade_8.png}
\caption{Goodness of Fit Plot for 2014 \textbf{\emph{Baseline}}
Referenced 8\(^{th}\) Grade ELA.}
\end{figure}

\subsubsection{SIMEX model fit}\label{simex-model-fit}

SIMEX model fit is also not expected to be perfect. In this case we
expect model misfit in the form of increased high SGPs for students with
lower prior performance (and a complementary decrease in low SGPs for
those students), and the reverse expectation for high achieving
students. This is visible in the goodness of fit plot in Figure 3 where
the SIMEX correction method has been applied to the 8\(^{th}\) grade ELA
baseline model.

\begin{figure}[htbp]
\centering
\includegraphics{img/Examples/Goodness_of_Fit/ELA.2014.BASELINE.SIMEX/gofSGP_Grade_8.png}
\caption{Goodness of Fit Plot for 2014 \textbf{\emph{SIMEX Corrected,}}
Baseline Referenced 8\(^{th}\) Grade ELA.}
\end{figure}

\subsection{Student Level Results}\label{student-level-results}

To investigate the possibility that individual level misfit might impact
summary level results, student growth percentile analyses were run on
all students and the results were examined relative to prior
achievement. With perfect fit to data, the correlation between students'
most recent prior achievement scores and their student growth
percentiles is zero (i.e., the goodness of fit tables would have a
uniform distribution of percentiles across all previous scale score
levels). To investigate in another way, correlations between prior
student scale scores and student growth percentiles were
calculated.\footnote{In addition to providing information about model
  fit, these student-level correlations can assess potential impact of
  test ceiling effects.}

Student-level correlations between the various SGP versions (uncorrected
and SIMEX corrected cohort and baseline-referenced estimates) and prior
achievement are presented here. The results are generally as expected.
With cohort-referenced percentiles, when the model is perfectly fit to
the data, the correlation between students' most recent prior
achievement scores and their student growth percentiles is zero (i.e.,
there is a uniform distribution of percentiles across all previous scale
score levels). Correlations for Georgia cohort-referenced SGPs are all
essentially zero. This provides assurance that the models have fit the
data well, and indicate that students can demonstrate high (or low)
growth regardless of prior achievement using cohort-referenced SGPs.

Baseline referenced growth percentiles relax assumptions about
correlations and the aspiration for a uniform distribution of
percentiles. Rather than considering a single year's cohort, baseline
SGPs are referenced against a ``super-cohort'' of several years of
students linked by common course/grade progressions. This allows us to
examine whether or not the system as a whole might be improving (or
declining) over time relative to the established baseline. That is, if
the system is improving over time, we would see that improvement in
median SGPs that are greater than 50 (more than half of students would
have growth greater than what \emph{was} typical growth in the past).
Furthermore, if students at different levels of prior achievement
experience differential growth a correlation greater or less than zero
may be observed.\footnote{A major assumption required in producing
  baseline-referenced SGPs is that the scale scores are well anchored.
  If this assumption does not hold, then any deviation from ``typical''
  growth may be purely and artifact of the test scaling procedure.
  However, it would seem that this type of ``scale drift'' would be
  consistent across scale score levels. Therefore it isn't clear that
  this would cause a differential impact on growth, and thereby a
  positive or negative correlation at the student level between growth
  and prior achievement.} For example, if lower achieving students had
consistently higher growth relative to the baseline cohort and higher
achieving students had consistently lower growth, then a negative
correlation would be expected at the student-level. Although small, many
correlations between baseline SGPs and prior achievement deviate from
zero.

SIMEX corrected SGPs induce a negative correlation between growth and
prior achievement (regardless of reference type) as shown in Table 9
below. Rather than a uniform distribution, SIMEX produces a distribution
in which growth for lower prior achieving students' growth is weighted
upward and higher prior achieving students' growth is weighted down. In
theory this produces biased student-level SGPs but may decrease the bias
in aggregate growth measures as has been documented previous SIMEX
reports. This in turn decreases the observed correlations between
``uncorrected'' SGPs and prior achievement\footnote{Note that
  aggregate-level correlations that are initially negative also become
  increasingly negative rather than return to zero. This has been noted
  in test results from other states not presented here.}.

\subsubsection{CRCT}\label{crct}

\begin{table}[!hbtp]
\begin{center}
\begin{normalsize}
\color{black}
\begin{tabular}{rrrrrr}
\multicolumn{6}{c}{\begin{minipage}[c]{0.764705882352941\linewidth}\caption*{\textbf{Table 9:} CRCT Correlations between Student-Level Prior Standardized Scale Score and SGP Versions}\end{minipage}}\\ 
\hline
Grade&Cohort SGP&Cohort SIMEX&Baseline SGP&Baseline SIMEX&N Size\\ 
\hline
\textbf{Reading}&&&&&\\ 
4&-0.04&-0.14&-0.05&-0.17&116,555\\ 
5&-0.01&-0.11& 0.04&-0.06&116,011\\ 
6& 0.00&-0.08&-0.15&-0.25&115,579\\ 
7& 0.01&-0.08& 0.05&-0.06&118,100\\ 
8& 0.00&-0.07& 0.23& 0.15&117,649\\ 
\textbf{ELA}&&&&&\\ 
4&-0.02&-0.12&-0.05&-0.17&116,512\\ 
5& 0.00&-0.09&-0.07&-0.16&115,884\\ 
6& 0.00&-0.08& 0.10& 0.01&115,387\\ 
7& 0.00&-0.07&-0.05&-0.16&117,865\\ 
8& 0.01&-0.07&-0.05&-0.14&117,354\\ 
\textbf{Mathematics}&&&&&\\ 
4&-0.02&-0.11&-0.04&-0.13&116,432\\ 
5& 0.00&-0.08& 0.04&-0.03&115,154\\ 
6& 0.00&-0.07&-0.06&-0.14&114,434\\ 
7& 0.01&-0.07& 0.09& 0.01&116,981\\ 
8& 0.01&-0.05& 0.21& 0.16&116,387\\ 
\textbf{Science}&&&&&\\ 
4&-0.01&-0.10&-0.05&-0.16&119,841\\ 
5& 0.00&-0.08&-0.10&-0.17&119,799\\ 
6& 0.00&-0.06& 0.01&-0.07&119,473\\ 
7& 0.00&-0.05& 0.00&-0.07&121,273\\ 
8& 0.00&-0.06&-0.08&-0.15&120,518\\ 
\textbf{Social Studies}&&&&&\\ 
4&-0.01&-0.10& 0.01&-0.08&119,417\\ 
5& 0.00&-0.07& 0.02&-0.06&119,367\\ 
6& 0.00&-0.06& 0.04&-0.02&119,045\\ 
7& 0.00&-0.06&-0.01&-0.08&120,844\\ 
8& 0.00&-0.06& 0.01&-0.06&119,989\\ 
\hline
\end{tabular}
\end{normalsize}
\end{center}
\end{table}

\color{black}

\subsubsection{EOCT Baseline Referenced
Subjects}\label{eoct-baseline-referenced-subjects}

\begin{table}[H]
\caption*{\textbf{Table 10:} Correlations between Student-Level Prior Standardized Scale Score and EOCT Baseline Referenced SGP Versions\label{table10}} 
\begin{center}
\begin{tabular}{lcllllcl}
\hline\hline
\multicolumn{1}{c}{\bfseries }&\multicolumn{1}{c}{\bfseries }&\multicolumn{4}{c}{\bfseries Aggregate Group Type}&\multicolumn{1}{c}{\bfseries }&\multicolumn{1}{c}{\bfseries }\tabularnewline
\cline{1-8}
\multicolumn{1}{c}{Content Area}&\multicolumn{1}{c}{}&\multicolumn{1}{c}{Cohort SGP}&\multicolumn{1}{c}{Cohort SIMEX}&\multicolumn{1}{c}{Baseline SGP}&\multicolumn{1}{c}{Baseline SIMEX}&\multicolumn{1}{c}{}&\multicolumn{1}{c}{N Size}\tabularnewline
\hline
Grade 9 Lit     && 0.00&-0.08&-0.19&-0.27&&121,063\tabularnewline
American Lit    && 0.01&-0.08&-0.02&-0.10&&106,673\tabularnewline
US History      &&-0.01&-0.09&-0.02&-0.10&&101,613\tabularnewline
Economics       && 0.00&-0.07&-0.03&-0.09&&95,447\tabularnewline
Biology         && 0.00&-0.08&-0.03&-0.10&&119,521\tabularnewline
Physical Science&& 0.00&-0.07&-0.02&-0.09&&78,867\tabularnewline
\hline
\end{tabular}\end{center}

\end{table}

\subsubsection{EOCT Cohort Referenced
Subjects}\label{eoct-cohort-referenced-subjects}

\begin{table}[!htbp]
\caption*{\textbf{Table 11:} Correlations between Student-Level Prior Standardized Scale Score and EOCT Cohort Referenced SGP Versions\label{table11}} 
\begin{center}
\begin{tabular}{lcllcl}
\hline\hline
\multicolumn{1}{c}{\bfseries }&\multicolumn{1}{c}{\bfseries }&\multicolumn{2}{c}{\bfseries Aggregate Group Type}&\multicolumn{1}{c}{\bfseries }&\multicolumn{1}{c}{\bfseries }\tabularnewline
\cline{1-6}
\multicolumn{1}{c}{Content Area}&\multicolumn{1}{c}{}&\multicolumn{1}{c}{Cohort SGP}&\multicolumn{1}{c}{Cohort SIMEX}&\multicolumn{1}{c}{}&\multicolumn{1}{c}{N Size}\tabularnewline
\hline
Coordinate Algebra&&0.00&-0.09&&129,534\tabularnewline
Analytic Geometry &&0.02&-0.08&&110,958\tabularnewline
Mathematics II    &&0.00&-0.08&& 7,805\tabularnewline
\hline
\end{tabular}\end{center}

\end{table}

\subsection{Group Level Results}\label{group-level-results}

Unlike when reporting SGPs at the individual level, when aggregating to
the group level (e.g., school) the correlation between aggregate prior
student achievement and aggregate growth is rarely zero. The correlation
between prior student achievement and growth at the school level is a
compelling descriptive statistic because it indicates whether students
attending schools serving higher achieving students grow faster (on
average) than those students attending schools serving lower achieving
students. Results from previous state analyses show a correlation
between prior achievement of students associated with a current school
(quantified as percent at/above proficient) and the median SGP to be
between 0.1 and 0.3. That is, these results indicate that on average,
students attending schools serving lower achieving students tend to
demonstrate less exemplary growth than those attending schools serving
higher achieving students. Equivalently, based upon ordinary least
squares (OLS) regression assumptions, the prior achievement level of
students attending a school accounts for between 1 and 10 percent of the
variability observed in student growth. There are no definitive numbers
on what this correlation should be, but studies on value-added models
show similar results (McCaffrey, Han, \& Lockwood, 2008).

\subsubsection{School Level Results}\label{school-level-results}

To illustrate these relationships visually, the bubble charts in Figures
4 through 8 depict growth as quantified by the median SGP of students at
the school against achievement/status, quantified by percentage of
student at/above proficient (advanced) at the school. The charts have
been successful in helping to motivate the discussion of the two
qualities: student achievement and student growth. Though the figures
are not detailed enough to indicate strength of relationship between
growth and achievement, they are suggestive and valuable for discussions
with stakeholders who are being introduced to the growth model for the
first time.

\begin{figure}[htbp]
\centering
\includegraphics{img/Bubble_Plots/2014/State/Style_1/Georgia_2014_ELA_State_Bubble_Plot_(Prior_Achievement).png}
\caption{School-level Bubble Plots for Georgia: ELA, 2013-2014.}
\end{figure}

\begin{figure}[htbp]
\centering
\includegraphics{img/Bubble_Plots/2014/State/Style_1/Georgia_2014_Reading_State_Bubble_Plot_(Prior_Achievement).png}
\caption{School-level Bubble Plots for Georgia: Reading, 2013-2014.}
\end{figure}

\begin{figure}[htbp]
\centering
\includegraphics{img/Bubble_Plots/2014/State/Style_1/Georgia_2014_Mathematics_State_Bubble_Plot_(Prior_Achievement).png}
\caption{School-level Bubble Plots for Georgia: Mathematics, 2013-2014.}
\end{figure}

\begin{figure}[htbp]
\centering
\includegraphics{img/Bubble_Plots/2014/State/Style_1/Georgia_2014_Science_State_Bubble_Plot_(Prior_Achievement).png}
\caption{School-level Bubble Plots for Georgia: Science, 2013-2014.}
\end{figure}

\begin{figure}[htbp]
\centering
\includegraphics{img/Bubble_Plots/2014/State/Style_1/Georgia_2014_Social Studies_State_Bubble_Plot_(Prior_Achievement).png}
\caption{School-level Bubble Plots for Georgia: Social Studies,
2013-2014.}
\end{figure}

The relationship between average prior student achievement and median
SGP observed for Georgia is relatively strong compared to some other
states for which the Center has done SGP analyses. Table 12 shows
correlations between prior achievement (measured as the mean prior
standardized scale score as well as the percent at/above proficient at
the school\footnote{Percent Prior Proficient in this case is determined
  by the percent of student's that scored in the Proficient or Advanced
  range of all student's that received a score. This measure does not
  reflect student's that did not receive a score but are included in the
  denominator of Percent Meeting Standard as displayed in the DOE
  Georgia Report Card.}). All results shown here are for schools with 15
or more students.

\begin{table}[H]
\caption*{\textbf{Table 12:} Correlations between Mean Prior Standardized Scale Score and Aggregate SGPs - (Combined Subjects)\label{table12}} 
\begin{center}
\begin{tabular}{lcllll}
\hline\hline
\multicolumn{1}{c}{\bfseries }&\multicolumn{1}{c}{\bfseries }&\multicolumn{4}{c}{\bfseries SGP Aggregate Type}\tabularnewline
\cline{1-6}
\multicolumn{1}{c}{Year}&\multicolumn{1}{c}{}&\multicolumn{1}{c}{Median SGP}&\multicolumn{1}{c}{Mean SGP}&\multicolumn{1}{c}{Median SIMEX}&\multicolumn{1}{c}{Mean SIMEX}\tabularnewline
\hline
2012&&0.60&0.61&0.45&0.47\tabularnewline
2013&&0.58&0.58&0.41&0.42\tabularnewline
2014&&0.54&0.55&0.38&0.38\tabularnewline
\hline
\end{tabular}\end{center}

\end{table}

Correlation tables describing the relationship between prior achievement
(defined here as mean prior standardized scale score) and aggregate
growth percentiles are presented below in separate subsections for CRCT
and EOCT subjects. The first correlation table in the each subsection
provides these overall SGP aggregates' relationships with mean prior
standardized scale scores. The additional correlation tables are
dis-aggregated by content area, and content area and grade to provide
more detail.

\textbf{\emph{CRCT}}

\begin{table}[H]
\caption*{\textbf{Table 13:} CRCT Correlations between School-Level Mean Prior Standardized Scale Score and (Baseline) SGP Aggregations.\label{table13}} 
\begin{center}
\begin{tabular}{lcllll}
\hline\hline
\multicolumn{1}{c}{\bfseries }&\multicolumn{1}{c}{\bfseries }&\multicolumn{4}{c}{\bfseries Aggregate SGP Type}\tabularnewline
\cline{1-6}
\multicolumn{1}{c}{Year}&\multicolumn{1}{c}{}&\multicolumn{1}{c}{Median SGP}&\multicolumn{1}{c}{Mean SGP}&\multicolumn{1}{c}{Median SIMEX}&\multicolumn{1}{c}{Mean SIMEX}\tabularnewline
\hline
&&&&&\tabularnewline
2012&&0.58&0.60&0.34&0.37\tabularnewline
2013&&0.58&0.60&0.33&0.34\tabularnewline
2014&&0.58&0.61&0.38&0.39\tabularnewline
\hline
&&&&&\tabularnewline
2012&&0.48&0.50&0.3&0.32\tabularnewline
2013&&0.48&0.49&0.28&0.28\tabularnewline
2014&&0.45&0.46&0.21&0.22\tabularnewline
\hline
&&&&&\tabularnewline
2012&&0.45&0.46&0.34&0.35\tabularnewline
2013&&0.46&0.48&0.34&0.36\tabularnewline
2014&&0.47&0.48&0.37&0.37\tabularnewline
\hline
&&&&&\tabularnewline
2012&&0.46&0.47&0.33&0.34\tabularnewline
2013&&0.31&0.32&0.12&0.14\tabularnewline
2014&&0.30&0.31&0.12&0.13\tabularnewline
\hline
&&&&&\tabularnewline
2012&&0.37&0.39&&\tabularnewline
2013&&0.29&0.30&0.15&0.17\tabularnewline
2014&&0.37&0.37&0.24&0.25\tabularnewline
\hline
\end{tabular}\end{center}

\end{table}

\begin{table}[H]
\caption*{\textbf{Table 14:} 2014 CRCT Correlations between School-Level prior standardized scale score and Aggregate (Baseline) SGPs by Grade.\label{table14}} 
\begin{center}
\begin{tabular}{lcllll}
\hline\hline
\multicolumn{1}{c}{\bfseries }&\multicolumn{1}{c}{\bfseries }&\multicolumn{4}{c}{\bfseries Aggregate Group Type}\tabularnewline
\cline{1-6}
\multicolumn{1}{c}{Grade}&\multicolumn{1}{c}{}&\multicolumn{1}{c}{Median SGP}&\multicolumn{1}{c}{Median SIMEX}&\multicolumn{1}{c}{Mean SGP}&\multicolumn{1}{c}{Mean SIMEX}\tabularnewline
\hline
&&&&&\tabularnewline
4&&0.31& 0.05&0.34& 0.05\tabularnewline
5&&0.51& 0.35&0.54& 0.37\tabularnewline
6&&0.15&-0.15&0.14&-0.14\tabularnewline
7&&0.54& 0.35&0.56& 0.38\tabularnewline
8&&0.76& 0.69&0.78& 0.72\tabularnewline
\hline
&&&&&\tabularnewline
4&&0.32& 0.10&0.33& 0.10\tabularnewline
5&&0.16&-0.02&0.15&-0.04\tabularnewline
6&&0.54& 0.42&0.55& 0.41\tabularnewline
7&&0.27& 0.03&0.30& 0.08\tabularnewline
8&&0.38& 0.20&0.41& 0.21\tabularnewline
\hline
&&&&&\tabularnewline
4&&0.26& 0.15&0.26& 0.15\tabularnewline
5&&0.23& 0.15&0.25& 0.17\tabularnewline
6&&0.23& 0.14&0.24& 0.14\tabularnewline
7&&0.53& 0.44&0.55& 0.45\tabularnewline
8&&0.54& 0.48&0.55& 0.50\tabularnewline
\hline
&&&&&\tabularnewline
4&&0.25& 0.07&0.27& 0.09\tabularnewline
5&&0.01&-0.11&0.03&-0.10\tabularnewline
6&&0.24& 0.14&0.25& 0.15\tabularnewline
7&&0.14& 0.03&0.15& 0.03\tabularnewline
8&&0.00&-0.11&0.03&-0.09\tabularnewline
\hline
&&&&&\tabularnewline
4&&0.22& 0.11&0.23& 0.12\tabularnewline
5&&0.05&-0.04&0.06&-0.03\tabularnewline
6&&0.16& 0.11&0.18& 0.12\tabularnewline
7&&0.17& 0.07&0.18& 0.09\tabularnewline
8&&0.04&-0.05&0.06&-0.03\tabularnewline
\hline
\end{tabular}\end{center}

\end{table}

\textbf{\emph{EOCT}}

\begin{table}[H]
\caption*{\textbf{Table 15:} 2011 to 2014 Correlations between Mean Prior Standardized Scale Score and Aggregate Baseline SGPs.\label{table15}} 
\begin{center}
\begin{tabular}{lcllll}
\hline\hline
\multicolumn{1}{c}{\bfseries }&\multicolumn{1}{c}{\bfseries }&\multicolumn{4}{c}{\bfseries Aggregate SGP Type}\tabularnewline
\cline{1-6}
\multicolumn{1}{c}{Year}&\multicolumn{1}{c}{}&\multicolumn{1}{c}{Median SGP}&\multicolumn{1}{c}{Mean SGP}&\multicolumn{1}{c}{Median SIMEX}&\multicolumn{1}{c}{Mean SIMEX}\tabularnewline
\hline
&&&&&\tabularnewline
2012&&0.50&0.53&0.37&0.4\tabularnewline
2013&&0.36&0.37&0.2&0.2\tabularnewline
2014&&0.02&0.04&-0.15&-0.16\tabularnewline
\hline
&&&&&\tabularnewline
2012&&0.48&0.50&0.32&0.35\tabularnewline
2013&&0.22&0.21&0.01&-0.01\tabularnewline
2014&&0.37&0.41&0.21&0.23\tabularnewline
\hline
&&&&&\tabularnewline
2012&&0.14&0.17&&\tabularnewline
2013&&0.07&0.08&-0.02&-0.01\tabularnewline
2014&&0.07&0.09&-0.03&-0.02\tabularnewline
\hline
&&&&&\tabularnewline
2012&&0.06&0.07&-0.04&-0.03\tabularnewline
2013&&0.19&0.21&0.11&0.12\tabularnewline
2014&&0.07&0.09&-0.01&0\tabularnewline
\hline
&&&&&\tabularnewline
2012&&0.16&0.17&0.03&0.05\tabularnewline
2013&&0.17&0.18&0.05&0.07\tabularnewline
2014&&0.13&0.14&0.02&0.03\tabularnewline
\hline
&&&&&\tabularnewline
2012&&0.23&0.24&0.13&0.14\tabularnewline
2013&&0.34&0.36&0.26&0.27\tabularnewline
2014&&0.29&0.29&0.19&0.21\tabularnewline
\hline
\end{tabular}\end{center}

\end{table}

\begin{table}[H]
\caption*{\textbf{Table 16:} 2011 to 2014 Correlations between Mean Prior Standardized Scale Score and Cohort SGPs.\label{table16}} 
\begin{center}
\begin{tabular}{lcllll}
\hline\hline
\multicolumn{1}{c}{\bfseries }&\multicolumn{1}{c}{\bfseries }&\multicolumn{4}{c}{\bfseries Aggregate SGP Type}\tabularnewline
\cline{1-6}
\multicolumn{1}{c}{Year}&\multicolumn{1}{c}{}&\multicolumn{1}{c}{Median SGP}&\multicolumn{1}{c}{Mean SGP}&\multicolumn{1}{c}{Median SIMEX}&\multicolumn{1}{c}{Mean SIMEX}\tabularnewline
\hline
&&&&&\tabularnewline
2013&&0.38&0.37&0.24&0.24\tabularnewline
2014&&0.40&0.40&0.25&0.24\tabularnewline
\hline
&&&&&\tabularnewline
2014&&0.47&0.48&0.32&0.33\tabularnewline
\hline
&&&&&\tabularnewline
2012&&0.39&0.40&&\tabularnewline
2013&&0.41&0.41&0.3&0.3\tabularnewline
2014&&0.36&0.38&0.27&0.3\tabularnewline
\hline
\end{tabular}\end{center}

\end{table}

\section*{References}\label{references}
\addcontentsline{toc}{section}{References}

Betebenner, D. W. (2009). Norm- and criterion-referenced student growth.
\emph{Educational Measurement: Issues and Practice}, \emph{28}(4),
42--51.

Betebenner, D. W., {VanIwaarden}, A., Domingue, B., \& Shang, Y. (2015).
\emph{SGP: An r package for the calculation and visualization of student
growth percentiles \& percentile growth trajectories.} Retrieved from
\url{https://github.com/CenterForAssessment/SGP}

Chernozhukov, V., Fern{á}ndez-Val, I., \& Galichon, A. (2010). Quantile
and probability curves without crossing. \emph{Econometrica},
\emph{78}(3), 1093--1125.

Koenker, R. (2005). \emph{Quantile regression}. Cambridge: Cambridge
University Press.

McCaffrey, D., Han, B., \& Lockwood, J. (2008). From data to bonuses: A
case study of the issues related to awarding teachers pay on the basis
of their students' progress. In \emph{National center for performance
incentives conference performance incentives: Their growing impact on
american k-12 education, vanderbilt university, nashville, tN}.

Shang, Y., VanIwaarden, A., \& Betebenner, D. W. (2015). Covariate
measurement error correction for student growth percentiles using the
sIMEX method. \emph{Educational Measurement: Issues and Practice},
\emph{34}(1), 4--14.
http://doi.org/\href{http://dx.doi.org/10.1111/emip.12058}{10.1111/emip.12058}

\bibliographystyle{plainnat}
\bibliography{/media/Data/Dropbox/Github_Repos/Packages/SGPreports/rmarkdown/univ_resources/educ.bib}




\end{document}


